\section{Threats to Validity}
\label{sec:threats}

\noindent\textbf{Construct validity.}
Our evaluation items are \emph{synthesized} from raw threads---made
self-contained and stripped of their resolution---which could dilute the
authenticity claim or leak the answer. Both risks are addressed by construction
(\S\ref{sec:normalization}): the verbatim original reporter text is released
with every item; the deterministic leak check flags any fix identifier
present in a question that the reporter had not stated before the
answer, and we report the residual leak rate of the released set; and every
synthesized item is independently reviewed by two researchers, with Cohen's
$\kappa$ reported over the full benchmark for the knowledge-type, leak, and
context-completeness judgments and every disagreement reconciled before
evaluation. The knowledge-type label itself is a construct: mislabelling a
contextual pair as parametric would admit an unanswerable question into the
no-context capability test. The deterministic hardening rule
(\S\ref{sec:normalization}) biases errors in the safe direction, and
condition-aware gradeability ensures a labelling error cannot
convert an unknowable identifier into a scored miss. Finally, gold answers are
thread-grounded: a terse maintainer pointer yields a terse gold. We do not
inflate such golds; instead the grading-time resolution oracle
(\S\ref{sec:grading}) supplies the resolved fix diff to the judge, and we
report results separately for pointer-style and substantive-answer items so the
two populations are not conflated.

\noindent\textbf{Internal validity.}
Free-text grading uses an LLM judge, and verifiable facts ground but do not eliminate
judge error---a response may cite the correct CVE while contradicting the
verdict it implies. We mitigate with the controls of \S\ref{sec:grading}:
per-claim verdicts with quoted evidence rather than holistic scores, a
$\ge 2$-judge ensemble with flagged disagreements, regex cross-checks of
identifier verdicts, and human calibration to $\kappa \ge 0.80$ on a stratified
sample before any reported run. The resolution oracle is scoped by a mechanical
eligibility rule (an explicit fix PR or commit reference, one resolution
attempt); items outside it are graded on the thread gold alone, and we report
eligibility and resolution rates so the oracle's coverage is visible. Snapshot
assembly is itself an instrument that can err: the snapshot commit is resolved
on the default branch (a question may concern a release branch), and the issue and advisory
corpora are bounded by what our miners retrieved. We release snapshot
identifiers (commit SHAs, corpus counts, database snapshot dates) with every
run for re-verification.

\noindent\textbf{External validity.}
We mine only publicly visible artifacts. Security issues are often disclosed
privately and resolved under embargo, leaving just an advisory or changelog in
public; SecDevQA captures questions negotiated in the open and may
under-represent privately routed ones, so our claims are scoped to publicly
resolved security Q\&A. Repositories are selected for catalogued security
activity and popularity (\dnum{three} languages so far), which may bias toward
projects with formal advisory processes; we report per-repository results and
cap any repository's share of the corpus to keep the bias bounded and visible.
The hard-verifiability filter excludes question types that rarely yield
checkable identifiers (e.g.\ severity assessment, design rationale); deriving
the taxonomy from the full corpus, including non-hard pairs, keeps that gap
measurable rather than silent.

\noindent\textbf{Contamination and the time cap.}
Public threads may be memorized~\cite{stackrepoqa2025}, and the time cap
constrains a system's \emph{context}, not its \emph{weights}: a model whose
training data includes the resolved thread can reproduce the resolution under
any condition. We therefore measure memorization rather than assume its
absence, with three instruments (\S\ref{sec:contamination}): no-context
accuracy on contextual items, training-cutoff stratification over $t_\theta$, and
transcript-level detection for the typed-tool agent (a fix identifier asserted
without the serving tool group ever being called). A residual limitation
remains for off-the-shelf agents: a shell-equipped agent cannot be provably
prevented from reaching the network, so the external-agent condition is
best-effort capped---the sandbox physically contains no post-report data and
the prompt forbids external access---and we release its transcripts; airtight
claims rest on the typed-tool reference agent.
